\begin{table}[h]
\centering
\small
\begin{tabular}{c|c|C{43pt}R{43pt}R{43pt}R{43pt}R{43pt}R{43pt}R{43pt}}
\toprule
\textbf{Dataset} &
& \textsc{Truth} & \indep & \markov & \hmm & \gan & \rbm & \method \\
\midrule
\multirow{3}{*}{\textbf{FIGS}}
& PCA1-2 & 0.0018 & 0.0856 & 0.0816 & 0.0390 & 0.0022 & \textbf{0.0020} & 0.0026 \\
& PCA3-4 & 0.0020 & 0.0857 & 0.0720 & 0.0151 & 0.0023 & \textbf{0.0020} & 0.0024 \\
& PCA5-6 & 0.0019 & 0.0658 & 0.0593 & 0.0273 & \textbf{0.0022} & 0.0022 & 0.0025 \\
\bottomrule
\end{tabular}
    \caption{\textbf{PCA distances in a common basis.} Wasserstein 2D (Sinkhorn, $\epsilon=2\times10^{-3}$) distances between the PCA representations of real (test set) versus generated individuals. A single PCA is fit once on the union of all eight datasets (training set, test set and the six sets of artificial genomes) and every dataset is projected into that same space, which is the space shown in Figures~\ref{fig:pca} and \ref{fig:pca-all}. Unlike Table~\ref{tab:global-distance}, where a separate PCA is fit per method, these values are directly comparable across methods. Truth represents the training set. Bolded values indicate the best among all compared models.}
\label{tab:pca-common-basis}
\end{table}
